\section{Geometric Properties of Embedding Spaces}

While the theoretical foundations of representation space geometry have been extensively formalised, prior empirical investigations into these properties have predominantly been confined to general-language domains. Foundational studies by \citet{ethayarajh-2019-contextual} and \citet{mu2018allbutthetop} exposed that contextualised embeddings from standard language models exhibit severe anisotropy, occupying a narrow cone rather than utilising the full representational capacity of the vector space. Building on these diagnostics, \citet{wang-isola-2020-contrastive} established alignment and uniformity as measurable quality metrics, proving their predictive value for downstream performance on standard Semantic Textual Similarity tasks.

To address these geometric pathologies, several post-hoc correction methods have emerged in the recent literature. While \citet{Li_2020_Sentence} introduced complex normalising flows, approaches such as the \ac{ABTT} transform \cite{mu2018allbutthetop} and Whitening transformations demonstrated that linear-algebraic interventions can achieve representational improvements \cite{su2021whitening}. However, these optimisation strategies were almost exclusively validated on general-domain sentence embeddings, leaving their efficacy in highly specialised, dense retrieval settings largely unexplored.
